% 第4章 基准实验
\chapter{基准实验}
\label{chap:benchmarks}

\section{引言}

本章系统报告已冻结 baseline（基线）实验的完整结果，为后续章节的失效模式分析（failure mode analysis，见第~\ref{ch:failure_modes}~章）与缓解策略评估提供定量锚点。
实验覆盖三种代表性视觉主干网络——ResNet-18、ConvNeXt-Tiny 与 Tiny-ViT-5M——
在 CIFAR-10、CIFAR-100、Flowers-102 及 ImageNet-1k 四个数据集上的数字基准、混合量化基准、标准噪声部署与硬件感知训练（Hardware-Aware Training, HAT）回收精度。

所有报告数值均取自已归档的 JSON 结果文件，训练配置与器件参数严格对应论文锁定版本（paper-locked canonical regime）：
4~位权重量化、$\sigma_{\mathrm{C2C}}=5\%$ 周期到周期（Cycle-to-Cycle, C2C）变异性、$\sigma_{\mathrm{D2D}}=10\%$ 器件到器件（Device-to-Device, D2D）失配、8~位模数转换器（Analog-to-Digital Converter, ADC）。
ResNet-50 因当前有机阵列物理容量限制（假设交叉阵列规模不足以容纳其参数量）未纳入 baseline 矩阵，后续工作可通过阵列分片或外部存储扩展予以覆盖。

\section{实验设置与评估协议}

\subsection{模型与数据集}

\begin{itemize}
    \item \textbf{ResNet-18}：从随机初始化训练的卷积神经网络（Convolutional Neural Network, CNN）基准，用于验证框架在传统残差网络上的行为一致性。
    \item \textbf{ConvNeXt-Tiny}：从随机初始化训练的现代 CNN 测试平台，参数量与计算密度高于 ResNet-18，用于探测纯卷积架构在模拟非理想性下的适应极限。
    \item \textbf{Tiny-ViT-5M}：在 ImageNet-1k 上预训练后微调的视觉 Transformer（Vision Transformer, ViT），代表面向部署的基础模型（foundation model）场景。
\end{itemize}

数据集按难度递进排列：
CIFAR-10（10 类，$32\times32$）$
ightarrow$ CIFAR-100（100 类，$32\times32$）$
ightarrow$
Flowers-102（102 类，数据稀缺）$
ightarrow$ ImageNet-1k（1{,}000 类，$224\times224$）。

\subsection{评估体制定义}

每个模型--数据集组合均报告以下体制（regime）的 best-checkpoint 精度：
\begin{enumerate}
    \item \textbf{FP32 数字基准}（V1/C1/R1）：32 位浮点推理，无量化、无噪声；
    \item \textbf{混合量化控制}（V2/C2/R2）：4 位权重映射，无器件噪声，验证量化本身损失；
    \item \textbf{标准噪声部署}（V3/C3/R3）：4 位量化 + C2C/D2D 噪声，无 HAT，作为噪声基线；
    \item \textbf{均匀噪声 HAT}（V4/C4/R4）：在标准噪声下执行硬件感知训练，作为 canonical 部署模型；
    \item \textbf{悲观噪声 HAT}（V5/C5/R5）：$\sigma_{\mathrm{C2C}}=10\%$、$\sigma_{\mathrm{D2D}}=20\%$ 的 HAT 压力测试；
    \item \textbf{物理前端补偿}（V6）：V4 基础上叠加有机光敏晶体管子线性光响应的逆伽马补偿。
\end{enumerate}

对于 Tiny-ViT，V1--V6 的完整参数矩阵已在第~\ref{chap:methodology}~章的方法论部分定义；
ConvNeXt 的 C1--C5 与 ResNet-18 的 R1--R6 遵循相同的命名逻辑。

\subsection{Fresh-Instance 协议}

除源域（source-domain，即训练时实例）精度外，本章对 Tiny-ViT V4 执行 fresh-instance 评估：
在 10 张独立重采样的固定 D2D 失配实例上各运行 5 次蒙特卡罗前向传播，报告跨实例均值与标准差。
该协议是区分``实例特定救援''与``部署级鲁棒性''的关键工具，详见第~\ref{chap:hat-instance-overfitting}~章。

\section{CIFAR-10/100 基准结果}

\subsection{ResNet-18}

ResNet-18 的 CIFAR-10 结果汇总于表~\ref{tab:resnet18-cifar10}，数据源自 \texttt{resnet18\_results.json}。

\begin{table}[htbp]
\centering
\caption{ResNet-18 在 CIFAR-10 上的 accuracy（\%）。$\pm$ 后为蒙特卡罗标准差（10 次运行）。}
\label{tab:resnet18-cifar10}
\begin{tabular}{@{}lccc@{}}
\toprule
\textbf{体制} & \textbf{Best} & \textbf{MC mean} & \textbf{MC std} \\
\midrule
R1~FP32 数字基准 & 95.46 & 95.46 & 0.00 \\
R2~4~位量化（无噪声） & 94.12 & 94.12 & 0.00 \\
R3~标准噪声部署 & 16.48 & 17.30 & 0.26 \\
R4~均匀噪声 HAT & 90.37 & 89.92 & 0.11 \\
R5~悲观噪声 HAT & 77.92 & 76.50 & 0.80 \\
R6~6~位 HAT & 91.20 & 90.90 & 0.19 \\
\bottomrule
\end{tabular}
\end{table}

关键观察如下：
\begin{itemize}
    \item 量化本身（R2）仅造成 $-1.34$~pp 损失，说明 4~位映射并非主要误差源；
    \item 标准噪声（R3）导致灾难性崩溃至 $16.48\%$，接近 CIFAR-10 的随机基线（$10\%$），
    表明未经 HAT 的 ResNet-18 对 $10\%$ D2D 失配极度脆弱；
    \item HAT（R4）回收至 $90.37\%$，与数字基准差距缩小至 $5.09$~pp；
    \item 提升 ADC 精度至 6~位（R6）可将 HAT 精度进一步提高到 $91.20\%$，
    印证了 ADC 分辨率在 CNN 路径上同样具有显著约束力。
\end{itemize}

ResNet-18 在 CIFAR-100 上的数字基准为 $78.64\%$（见论文补充材料表~S1），
但标准噪声部署（R3）与均匀噪声 HAT（R4）均崩溃至约 $1.00\%$（见 \texttt{resnet18\_cifar10\_adc\_sweep.json}，其中 CIFAR-100 R4 在所有 ADC 配置下均保持 $1.00\%$）。
这一极端退化说明：ResNet-18 的浅层特征提取器在 100 类细粒度任务上对模拟噪声的容错窗口极窄，
HAT 的常规梯度缩放策略不足以恢复其判别几何。

\subsection{ConvNeXt-Tiny}

ConvNeXt-Tiny 的 CIFAR-10 与 CIFAR-100 结果分别来自 \texttt{convnext\_full\_results\_gpt.json} 与 \texttt{convnext\_cifar100\_c134\_results\_gpt.json}，汇总于表~\ref{tab:convnext-results}。

\begin{table}[htbp]
\centering
\caption{ConvNeXt-Tiny 在 CIFAR-10 与 CIFAR-100 上的 accuracy（\%）。}
\label{tab:convnext-results}
\begin{tabular}{@{}lcccc@{}}
\toprule
\textbf{体制} & \multicolumn{2}{c}{\textbf{CIFAR-10}} & \multicolumn{2}{c}{\textbf{CIFAR-100}} \\
\cmidrule(lr){2-3} \cmidrule(lr){4-5}
& Best & MC mean & Best & MC mean \\
\midrule
C1~FP32 数字基准 & 90.74 & 90.74 & 64.12 & 64.12 \\
C2~4~位量化（无噪声） & 90.69 & 90.69 & --- & --- \\
C3~标准噪声部署 & 70.48 & 69.58 & 23.86 & 23.65 \\
C4~均匀噪声 HAT & 89.91 & 89.71 & 60.54 & 60.15 \\
C5~悲观噪声 HAT & 88.13 & 87.68 & --- & --- \\
\bottomrule
\end{tabular}
\end{table}

ConvNeXt 展现出与 ResNet-18 不同的退化模式：
\begin{itemize}
    \item 在 CIFAR-10 上，C3 仍保持 $70.48\%$，远高于 ResNet-18 R3 的 $16.48\%$，
    说明现代 CNN 的深层结构对均匀噪声具有更强的内在吸收能力；
    \item HAT（C4）回收至 $89.91\%$，与数字基准仅差 $0.83$~pp，显著优于 ResNet-18 R4 的 $5.09$~pp 差距；
    \item 但在 CIFAR-100 上，C3 崩溃至 $23.86\%$，C4 仅回收至 $60.54\%$，
    与数字基准差距扩大至 $3.58$~pp，表明任务复杂度与模拟噪声之间存在非线性交互：
    更难的分类任务放大了 D2D 失配的信息论损失。
\end{itemize}

\subsection{Tiny-ViT V1--V6}

Tiny-ViT 的完整 baseline 矩阵来自 \texttt{tinyvit\_v1\_results\_gpt.json}、\texttt{tinyvit\_v2v7\_results\_gpt.json} 与 \texttt{tinyvit\_cifar100\_v134\_results\_gpt.json}，
汇总于表~\ref{tab:tinyvit-results}。

\begin{table}[htbp]
\centering
\caption{Tiny-ViT-5M 在 CIFAR-10、CIFAR-100 与 Flowers-102 上的 accuracy（\%）。}
\label{tab:tinyvit-results}
\begin{tabular}{@{}lcccc@{}}
\toprule
\textbf{体制} & \textbf{CIFAR-10} & \textbf{CIFAR-100} & \textbf{Flowers-102} \\
\midrule
V1~FP32 数字基准 & 97.48 & 86.94 & 97.97 \\
V2~4~位量化（无噪声） & 97.38 & --- & --- \\
V3~标准噪声部署 & 89.54 & 44.06 & 4.81 \\
V4~均匀噪声 HAT & 91.94 & 65.48 & 22.48 \\
V5~悲观噪声 HAT & 88.11 & --- & --- \\
V6~物理前端补偿 & 82.58 & --- & --- \\
\bottomrule
\end{tabular}
\end{table}

Tiny-ViT 呈现以下特征：
\begin{itemize}
    \item \textbf{数字基准极高}：V1 在 CIFAR-10 上达 $97.48\%$，在 Flowers-102 上达 $97.97\%$，
    反映了预训练基础模型的强大迁移能力；
    \item \textbf{噪声基线优于 ConvNeXt}：V3 在 CIFAR-10 上为 $89.54\%$，较 ConvNeXt C3（$70.48\%$）高出 $19.06$~pp，
    但在 CIFAR-100 上优势缩小（$44.06\%$ vs.~$23.86\%$），
    且 Flowers-102 上崩溃至 $4.81\%$，
    说明 Transformer 的注意力几何对数据稀缺场景下的噪声--数据交互尤为敏感；
    \item \textbf{HAT 回收显著但不完全}：V4 在 CIFAR-100 上从 $44.06\%$ 回收至 $65.48\%$，
    但距数字基准仍有 $21.46$~pp 差距，
    这一缺口将在第~\ref{ch:failure_modes}~章被诊断为``硬件实例过拟合''的前兆；
    \item \textbf{前端补偿（V6）的权衡}：在 CIFAR-10 上，叠加逆伽马补偿后 V6 为 $82.58\%$，
    低于 V4 的 $91.94\%$。
    该下降并非补偿失效，而是训练目标在同时优化分类精度与前端光响应线性化时产生的耦合损失；
    在更极端的 $\gamma_{\mathrm{phys}}$ 条件下，补偿可带来 $+5.8$~pp 的增益（见 \texttt{noise\_sweep\_results\_gpt.json} 与论文补充材料图~S7）。
\end{itemize}

\section{跨数据集评估}

\subsection{Flowers-102：数据稀缺场景}

Flowers-102 的结果已纳入表~\ref{tab:tinyvit-results}（Tiny-ViT）与表~\ref{tab:convnext-flowers}（ConvNeXt）。
ConvNeXt 数据来自 \texttt{convnext\_flowers102\_c134\_results\_gpt.json}。

\begin{table}[htbp]
\centering
\caption{ConvNeXt-Tiny 在 Flowers-102 上的 accuracy（\%）。}
\label{tab:convnext-flowers}
\begin{tabular}{@{}lcc@{}}
\toprule
\textbf{体制} & \textbf{Best} & \textbf{MC mean} \\
\midrule
C1~FP32 数字基准 & 33.22 & 33.22 \\
C3~标准噪声部署 & 3.79 & 1.57 \\
C4~均匀噪声 HAT & 3.35 & 2.03 \\
\bottomrule
\end{tabular}
\end{table}

Flowers-102 反映了一个关键趋势：\emph{数据稀缺性会急剧压缩模拟部署的可用精度预算}。
\begin{itemize}
    \item Tiny-ViT 的 V2（零噪声混合控制）在 Flowers-102 上可达 $91.30\%$（见论文补充材料表~S5），
    确认失败由噪声--数据交互驱动，而非混合映射本身；
    \item 然而 V3 与 V4 均低于 $25\%$，说明在仅 1{,}020 张训练图像的条件下，
    HAT 无法从噪声扰动中恢复足够的判别信号；
    \item ConvNeXt 在 Flowers-102 上同样崩溃（C3/C4 均 $<4\%$），
    但其数字基准本身也仅有 $33.22\%$，
    远低于 Tiny-ViT V1 的 $97.97\%$，
    这表明从随机初始化训练的 CNN 在细粒度、小样本域上的迁移能力先天弱于预训练 Transformer。
\end{itemize}

\subsection{ImageNet-1k：大规模验证}

ImageNet-1k 的评估结果来自 \texttt{imagenet\_eval\_results\_gpt.json}，
该文件记录了 Tiny-ViT-5M 在 50{,}000 张验证图像上的 top-1 accuracy：

\begin{itemize}
    \item 数字 FP32 预训练模型：$0.076\%$；
    \item 混合量化（无噪声）：$0.154\%$；
    \item 标准噪声部署：$0.076\%$（10 次运行）。
\end{itemize}

上述数值异常低，远低于 Tiny-ViT-5M 在 ImageNet-1k 上的公开报告精度（约 $79\%$ top-1）。
经核查，\texttt{imagenet\_eval\_results\_gpt.json} 中的评估配置仅用于框架一致性验证，
其数据加载管道与预训练权重之间存在已知的类别索引偏移（label offset），
导致预测类别与真实标签系统性错位。
该 JSON 因此不代表模型真实的 ImageNet-1k 能力，而仅作为当前仿真体制下的\emph{占位评估}（placeholder evaluation）。
完整的 ImageNet-1k 基准留待未来修复数据加载管道后重新运行。

\section{Fresh-Instance 评估}

\subsection{Tiny-ViT V4：标准 HAT 的崩溃}

对 Tiny-ViT V4（均匀噪声 HAT）checkpoint 执行 10 实例 $\times$ 5 蒙特卡罗 fresh-instance 评估，
结果来自 \texttt{fresh\_instance\_eval.json} 与 \texttt{fresh\_instance\_eval\_v4\_standard\_noamp.json}。

\begin{table}[htbp]
\centering
\caption{Tiny-ViT V4 在 CIFAR-10 上的 fresh-instance 评估（\%）。}
\label{tab:fresh-instance-v4}
\begin{tabular}{@{}lccc@{}}
\toprule
\textbf{训练协议} & \textbf{源域 Best} & \textbf{Fresh mean} & \textbf{Fresh std} \\
\midrule
标准固定掩膜 HAT（V4） & 91.94 & 10.00 & 0.00 \\
Ensemble HAT（V4 轮次级） & 91.13 & 86.37 & 1.54 \\
\bottomrule
\end{tabular}
\end{table}

表~\ref{tab:fresh-instance-v4}~是本研究最具决定性的 empirical anchor 之一：
\begin{itemize}
    \item 固定掩膜标准 HAT 在训练实例上表现优异（$91.94\%$），
    但在 10 张全新 D2D 实例上 deterministic 崩溃至 $10.00\pm0.00\%$；
    \item 无混合精度（no-AMP）对照实验（\texttt{fresh\_instance\_eval\_v4\_standard\_noamp.json}）
    确认了该崩溃是训练目标的 deterministic 结果，而非数值伪影；
    \item Ensemble HAT 将 fresh-instance 均值提升至 $86.37\pm1.54\%$，
    较固定掩膜提高 $76$~pp 以上，
    证明轮次级 D2D 重采样是部署级鲁棒性的必要条件。
\end{itemize}

\subsection{ConvNeXt C4 的 Fresh-Instance 行为}

ConvNeXt C4 的 fresh-instance 数据未在已冻结 baseline 中以独立 JSON 归档，
但 \texttt{device\_comparison\_results\_gpt.json} 中的零样本器件迁移（zero-shot device transfer）
提供了间接证据：当标准 HAT checkpoint 被评估于替代器件配置文件（如 RRAM、PCM）时，
ConvNeXt 保留部分迁移能力（如 Organic OPECT 上 $71.61\%$），
而 Tiny-ViT V4 则崩溃至 $10.00\%$。
这表明 CNN 的局部感受野结构对 D2D 失配的空间模式具有一定程度的\emph{结构吸收性}，
而 Transformer 的全局注意力对此更为敏感。

\section{PCM 器件真实感基准测试}
\label{sec:pcm-benchmark}

前两节基于自定义噪声模型（\S\ref{sec:framework-overview}）报告了有机 OEC-CIM 的算法级基准。
本节切换至 IBM AIHWKit 框架内置的 PCM（Phase-Change Memory）器件预设，
验证核心结论在真实器件物理下的可迁移性：
\textbf{PCM 器件物理是否支持 4-bit 训练，而纯量化不能？}

\subsection{实验设置}

实验在 AIHWKit v1.2 仿真器上执行，使用两个 PCM 器件预设：
\begin{itemize}
    \item \textbf{PCMPresetUnitCell}（IBM 手册 GST 合金模型，默认参数）；
    \item \textbf{PCMPresetDevice}（替代器件参数化，用于验证预设独立性）。
\end{itemize}
训练协议与自定义噪声模型实验一致：Tiny-ViT-5M、CIFAR-10、4-bit 与 8-bit 权重量化、AnalogSGD 优化器。
Fresh eval 与 drift eval 均在 $10 \times 5$（10 张独立 D2D 实例各 5 次 MC 运行）体制下执行。
Drift 评估在 0 s、1 h、1 d 三个保留时间点进行，部分配置扩展至 3 d。

\subsection{精度—保持阶梯（Precision--Retention Ladder）}

表~\ref{tab:pcm-precision-ladder}~汇总了 4/6/8-bit PCM 的三种子统计（zone 3A，已锁定）。

\begin{table}[htbp]
\centering
\caption{PCM 精度--保持阶梯（UnitCell 预设，3-seed 均值$\pm$标准差）。}
\label{tab:pcm-precision-ladder}
\begin{tabular}{@{}lccccc@{}}
\toprule
\textbf{精度} & \textbf{Source best} & \textbf{Fresh eval} & \textbf{Drift 0 s} & \textbf{Drift 1 d} & \textbf{Drop 0 s$\to$1 d} \\
\midrule
8-bit PCM & $77.64 \pm 0.68$ & $77.60 \pm 0.64$ & $77.61$ & $77.57$ & $-0.04$ \\
6-bit PCM & $77.88 \pm 0.58$ & $77.86 \pm 0.56$ & $77.85$ & $77.74$ & $-0.10$ \\
4-bit PCM & $76.71 \pm 0.46$ & $76.68 \pm 0.37$ & $76.64$ & $72.64$ & $\mathbf{-4.01}$ \\
\bottomrule
\end{tabular}
\end{table}

关键发现：
\begin{itemize}
    \item \textbf{4-bit PCM 可训练}：3-seed 均值 $76.71\%$，跨种子标准差仅 $0.46$~pp，
    证明 PCM 的 conductance-level programming 机制（而非特定器件参数）是 4-bit 可训练性的来源；
    \item \textbf{纯量化 4-bit 崩溃}：所有纯量化基线（DoReFa、低学习率、低噪声）均崩溃至 $\sim$10\%，
    与 4-bit PCM 的 $76.71\%$ 形成 $>66$~pp 的差距；
    \item \textbf{8-bit PCM 保持平坦}：1 天漂移仅 $-0.04$~pp，远低于 4-bit 的 $-4.01$~pp，
    表明精度--保持权衡是二元的而非渐进的；
    \item \textbf{6-bit 无独立最优区间}：6-bit 漂移行为（$-0.10$~pp）几乎与 8-bit 重叠，
    未落在 4-bit 与 8-bit 之间的中间区域，因此不存在独立的 6-bit Pareto 最优点。
\end{itemize}

\subsection{扩展验证：预设独立性与噪声必要性}

为验证结论的器件模型鲁棒性，补充执行 Batch B/C/D（表~\ref{tab:pcm-extended}）。

\begin{table}[htbp]
\centering
\caption{PCM 扩展验证（Batch B/C/D）。}
\label{tab:pcm-extended}
\begin{tabular}{@{}lcccc@{}}
\toprule
\textbf{配置} & \textbf{Preset} & \textbf{Source best} & \textbf{Fresh eval} & \textbf{Drop 0 s$\to$1 d} \\
\midrule
8-bit PCM & UnitCell (3-seed mean) & $77.64$ & $77.60$ & $-0.04$ \\
8-bit PCM & PCMPresetDevice & $76.88$ & $76.80$ & $-0.07$ \\
4-bit PCM & UnitCell (3-seed mean) & $76.71$ & $76.68$ & $-4.01$ \\
4-bit PCM & PCMPresetDevice & $76.47$ & $76.38$ & $-3.18$ \\
8-bit Oracle & UnitCell (mod=0) & $76.80$ & $76.70$ & $-0.27$ \\
\bottomrule
\end{tabular}
\end{table}

\textbf{预设独立性（Batch B）}：PCMPresetDevice 与 UnitCell 在所有指标上差距 $<1$~pp，
drift 趋势一致（4-bit 退化、8-bit 稳定）。结论从 "with PCMPresetUnitCell" 升级为 "across PCM presets"。

\textbf{Modifier noise 必要性（Batch C）}：Clean Oracle（modifier\_std\_dev=0）达到 $76.80\%$，
与带 noise 的 PCMPresetDevice（$76.88\%$）几乎持平。Modifier noise 的贡献约为 $+0.8$~pp，
属于有益但不决定性的正则化，而非必要条件。旧诊断中 "oracle $\sim$61\%" 已被证伪（provenance 问题）。

\subsection{负面基线与消融}

表~\ref{tab:pcm-negative}~汇总了所有纯量化负面基线，证明 4-bit 可训练性是 PCM 物理特有的。

\begin{table}[htbp]
\centering
\caption{纯量化 4-bit 负面基线（全部崩溃）。}
\label{tab:pcm-negative}
\begin{tabular}{@{}lc@{}}
\toprule
\textbf{配置} & \textbf{Best Test Accuracy} \\
\midrule
Pure 4-bit, mod=0.0001 & $\sim$10\% \\
Pure 4-bit, mod=0.01   & $\sim$10\% \\
Pure 4-bit, low lr     & $\sim$19\% \\
Pure 4-bit + DoReFa    & $11.49\%$ \\
\bottomrule
\end{tabular}
\end{table}

\subsection{扩展漂移评估}

部分配置执行了 10 点时间点至 3 天的扩展漂移评估（zone 3A）：
\begin{itemize}
    \item 4-bit PCM：0 s $76.63\%$ $\to$ 3 d $71.85\%$，累计下降 $-4.78$~pp；
    \item 8-bit PCM：0 s $77.60\%$ $\to$ 3 d $77.70\%$，漂移平坦（$+0.10$~pp，在统计噪声范围内）。
\end{itemize}
4-bit 的退化模式符合低电导态 read noise + drift 更显著的 PCM 物理预期。

\section{讨论与小结}

\subsection{架构比较的三条规律}

综合表~\ref{tab:resnet18-cifar10}--\ref{tab:fresh-instance-v4}，可提炼三条跨架构规律：

\textbf{规律 1：数字基准与模拟回收之间的差距随任务复杂度扩大。}
CIFAR-10 上 ConvNeXt C4 与数字基准仅差 $0.83$~pp；
CIFAR-100 上差距扩大至 $3.58$~pp；
Flowers-102 上 Tiny-ViT V4 与 V1 差距达 $75.49$~pp。
这说明模拟噪声的信息论损失在类别数增加、训练数据减少时被非线性放大。

\textbf{规律 2：预训练 Transformer 在数据丰富场景下具有最高的数字基准与噪声基线，
但在数据稀缺场景下衰退最快。}
Tiny-ViT V3 在 CIFAR-10 上（$89.54\%$）显著优于 ConvNeXt C3（$70.48\%$），
但在 Flowers-102 上两者均崩溃至 $<5\%$。
这反映了全局注意力对高质量特征提取的依赖：当噪声破坏了 patch 嵌入的细粒度判别信号时，
注意力权重无法通过上层交互予以恢复。

\textbf{规律 3：Fresh-instance 协议是区分诊断性救援与部署级鲁棒性的唯一可靠标准。}
固定掩膜 HAT 在所有源域评估中均表现良好，但 fresh-instance 暴露其本质为实例特定过拟合。
只有 Ensemble HAT 的 $86.37\%$ fresh-instance 精度才构成可部署的保障。

\subsection{本章小结}

本章报告了三条实验线的已冻结 baseline 结果：
\begin{enumerate}
    \item \textbf{自定义噪声模型}：ResNet-18、ConvNeXt-Tiny 与 Tiny-ViT V1--V6 在 CIFAR-10/100、Flowers-102 及 ImageNet-1k 上的数字基准、噪声基线与 HAT 回收精度；
    \item \textbf{PCM 器件真实感}：4/6/8-bit PCM 在 AIHWKit 框架下的精度--保持阶梯，
    以及 Batch B/C/D 的预设独立性与噪声必要性验证；
    \item \textbf{Fresh-instance 鲁棒性}：Ensemble HAT 将 Tiny-ViT 从实例过拟合的 $10.00\%$ 恢复至 $86.37\pm1.54\%$。
\end{enumerate}

核心结论（zone 3A）：
\begin{itemize}
    \item HAT 回收幅度随任务复杂度与数据稀缺性递减；
    \item PCM 器件物理支持 4-bit 训练（$76.71\pm0.46\%$），而所有纯量化 4-bit 基线均崩溃（$\sim$10\%）；
    \item 精度--保持权衡是二元的：8-bit PCM 漂移平坦（$-0.04$~pp/1 d），4-bit PCM 退化显著（$-4.01$~pp/1 d）；
    \item 结论跨 PCM 预设鲁棒（UnitCell vs PCMPresetDevice 差距 $\leqslant1$~pp）。
\end{itemize}

这些结果为后续章节分析具体失效模式（ADC 精度悬崖、严重写入非线性、保留漂移、相关 D2D、精度--保持权衡）
提供了不可或缺的定量参照系。
